\subsection{MULTIPLICATION IN R}

The topology of the rational line \(\mathbf{Q}\) is compatible not only with the additive group structure, but also with the field structure of \(\mathbf{Q}\) . For the function \(xy\) is continuous at \((0, 0) \in \mathbf{Q} \times \mathbf{Q}\) , since for each integer \(n > 0\) the relations \(|x| \leqslant 1/n\) and \(|y| \leqslant 1/n\) together imply that \(|xy| \leqslant 1/n^2 \leqslant 1/n\) ; on the other hand, if \(a\) is any non-zero rational number, the function \(ax\) is continuous at \(x = 0\) , since for each integer \(n > 0\) the relation \(|x| \leqslant 1/n |a|\) implies that \(|ax| \leqslant 1/n\) . This shows that \(xy\) is continuous at every point of \(\mathbf{Q} \times \mathbf{Q}\) (Chapter III, § 6, no. 3).

To show that \(1 / x\) is continuous on \(\mathbf{Q}^*\) we shall establish more precisely that \(1 / x\) is uniformly continuous (with respect to the additive structure) in the complement of any neighbourhood \(\mathbf{V}\) of 0. Namely we have \(\left|\frac{1}{x} - \frac{1}{y}\right| = \frac{|x - y|}{xy}\) ; there exists an integer \(m > 0\) such that \(|x| \geqslant 1 / m\) for each \(x \in \mathbb{C}\mathbf{V}\) ; if \(x\) and \(y\) are any two points of \(\mathbb{C}\mathbf{V}\) such that \(|x - y| \leqslant 1 / m^2 n\) , we shall then have \(\left|\frac{1}{x} - \frac{1}{y}\right| \leqslant \frac{1}{n}\) .

The image, under the function \(1 / x\) , of any Cauchy filter on \(\mathbf{Q}^*\) (with respect to the additive uniformity) which does not have \(0\) as a cluster point, is a Cauchy filter (with respect to the additive uniformity). Hence (Chapter III, § 6, no. 8, Proposition 7):

PROPOSITION 1. The functions xy and 1/x, defined respectively on Q × Q and Q*, can be extended by continuity to R × R and R* respectively, and define a field structure on R. Endowed with this structure, R is called the field of real numbers.

All the properties of topological fields established in § 6 of Chapter III are of course applicable; in particular, every rational function of n real variables, with real coefficients, is continuous at every point of \(R^{n}\) where its denominator does not vanish.

